\section{Related Work}
\label{sec:related}

\paragraph{Traditional and learned video codecs.}
Standardized hybrid codecs HEVC~\cite{sullivan2012overview} and
VVC~\cite{bross2021overview} structure a sequence into intra (I), predicted
(P), and bidirectional (B) frames, transmitting residuals against motion-
compensated predictions through DCT and entropy coding. Learned video
codecs largely retain this GOP layout but replace each component with
neural modules: DCVC and its
descendants~\cite{li2024neural,jia2025dcvcrt} use learned conditional
entropy models and feature-space prediction. These methods now match or
exceed VVC in PSNR at the standard 0.05--0.20~bpp operating range with
real-time-class decode speeds, but their RD curves rapidly flatten below
$0.01$~bpp.

\paragraph{Generative video compression.}
A separate line of work uses generative models to sample a plausible
reconstruction at ultra-low bitrates rather than regress to a deterministic
prediction. GLC-Video~\cite{qi2025generative} and
GNVC-VD~\cite{mao2025gnvcvd} train end-to-end generative codecs that
achieve $0.10$--$0.15$ LPIPS at $0.003$~bpp, well below the LPIPS of
traditional and learned codecs at the same rate, but require dedicated
training. Free-GVC~\cite{ling2026freegvc} avoids training by performing
reverse channel coding along the trajectory of a pretrained CogVideoX
generator. GVCC~\cite{zeng2026gvcc} is the closest predecessor to our
work: it converts a pretrained rectified-flow video generator into a
stochastic process and transmits codebook-indexed per-step noise. We
inherit GVCC's codec backbone---including its Turbo-DDCM~\cite{vaisman2025turboddcm}
multi-atom codebook construction---and focus exclusively on closing the
decode-time gap left open by that work.

\paragraph{Multi-resolution diffusion sampling.}
The idea of running the early high-noise iterations of a diffusion or flow
model at lower resolution is well established in the sampling literature.
Spectral Progressive Diffusion (SPD)~\cite{xiao2026spd} provides a clean
spectral-domain reading: at each timestep $t$, only those frequency bands
whose signal-to-noise ratio crosses a threshold need to be active, and
lower bands cross first. Cascaded
diffusion~\cite{ho2022cascaded} and several text-to-image
systems~\cite{podell2023sdxl,esser2024sd3} exploit the same structure with
multiple separately trained models.

Our use of multi-resolution coding differs from these in two ways. First,
the goal is not to generate plausible samples but to drive an encoder-
specified target with bit-exact reproducibility on both sides. Second, the
low and high stages share a single pretrained backbone (no cascading of
distinct models), which forces a precise specification of \emph{what} the
low stage's target latent is---a question that does not arise in
generation, where any plausible low-frequency content is acceptable. The
target-domain mismatch we identify in Section~\ref{sec:method-mr} arises
exactly because the low stage is anchored by the encoder, not free to drift.

\paragraph{Accelerating diffusion sampling.}
Reducing the number of model evaluations is well studied in the
distillation literature: progressive distillation~\cite{salimans2022progressive},
consistency models~\cite{song2023consistency,luo2023latent}, and
Sana~\cite{xie2024sana}-style mobile distillations all aim to bring sample
counts from $T = 20$--$50$ down to $T \le 4$. These approaches require
retraining and produce a distilled generator that may not be compatible
with the encoder's per-step codebook reproducibility requirement. Closer
to us are inference-time caching schemes that reuse intermediate
quantities across denoising steps: DeepCache~\cite{ma2024deepcache} caches
internal U-Net features, and several recent
works~\cite{kim2025inferencetime} cache velocity outputs directly. We
adopt the latter family but cache the clean-image \emph{endpoint} rather
than the velocity (Section~\ref{sec:method-skip}); this choice is forced by
the codec requirement that the cache stay valid across an SDE step where
the trajectory state changes by an encoder-selected codebook innovation.

\paragraph{Pretrained video generators as the decoder backbone.}
We instantiate GVCCTurbo on Wan~2.1~\cite{wan2025}, a rectified-flow
video generator supporting unconditional (T2V), single-anchor (I2V), and
dual-anchor (FLF2V) generation modes at $480$p and $720$p+. The choice
follows~\cite{zeng2026gvcc} for direct comparability; the framework
inherits whatever resolution and conditioning support the base generator
provides.
